% MLIR-style composition of arrus.tx_rx across two dialect levels.
%
% Illustrates that an operation can contain other operations (MLIR
% regions-inside-ops pattern). The high-level `arrus.tx_rx` (Level 1,
% arrus dialect) is composed of sub-ops that compute delays from SoS,
% compose a TX descriptor, and acquire. The `arrus.acquire` sub-op is
% itself lowered by `Us4RBackend.compile` into a body of Level-2 ops
% from the `us4r` dialect (delays -> device cycles -> set_delays_register
% -> ... -> fire).
%
% The nesting is real MLIR-style: a composite op owns a region containing
% its body. Level 2 ops live INSIDE Level 1 ops after lowering; the graph
% is not two separate layers but one graph with nested op scopes.
%
% Build:  pdflatex ndarray_fig4_mlir_composition.tex

\documentclass[tikz,border=6pt]{standalone}
% Shared preamble for NdArray graph figures.
% Include with: % Shared preamble for NdArray graph figures.
% Include with: % Shared preamble for NdArray graph figures.
% Include with: \input{ndarray_common}
%
% Provides a consistent color palette and node styles for the three figures
% in docs/td/ndarray_design.md.

\usepackage{tikz}
\usetikzlibrary{positioning, arrows.meta, fit, calc, backgrounds, shapes.geometric}

% --- Device / role color palette ------------------------------------------------
\definecolor{arrusUs4R}   {HTML}{2CA9BC}   % ultrasound device (teal)
\definecolor{arrusGPU}    {HTML}{F28C28}   % CUDA GPU (orange)
\definecolor{arrusCPU}    {HTML}{4CAF50}   % CPU (green)
\definecolor{arrusHost}   {HTML}{9E9E9E}   % host / control (gray)
\definecolor{arrusBuffer} {HTML}{616161}   % stream buffer (dark gray)
\definecolor{arrusVariable}{HTML}{F4C542}  % Variable (gold)
\definecolor{arrusTrigger}{HTML}{D64545}   % Trigger source (red)
\definecolor{arrusOutput} {HTML}{5C6BC0}   % Graph output / sink (indigo)

% --- Node styles ---------------------------------------------------------------
\tikzset{
  op/.style       = {rectangle, rounded corners=2pt, draw=black!60, thick,
                     minimum width=26mm, minimum height=8mm, align=center,
                     font=\small},
  op/us4r/.style  = {op, fill=arrusUs4R!25,   draw=arrusUs4R!80!black},
  op/gpu/.style   = {op, fill=arrusGPU!25,    draw=arrusGPU!80!black},
  op/cpu/.style   = {op, fill=arrusCPU!25,    draw=arrusCPU!70!black},
  op/host/.style  = {op, fill=arrusHost!20,   draw=arrusHost!70!black},
  op/output/.style= {op, fill=arrusOutput!25, draw=arrusOutput!80!black},
  op/variable/.style = {ellipse, draw=arrusVariable!70!black, thick,
                        fill=arrusVariable!30, minimum width=22mm,
                        minimum height=7mm, align=center, font=\small\itshape},
  op/trigger/.style  = {regular polygon, regular polygon sides=3,
                        draw=arrusTrigger!80!black, thick, fill=arrusTrigger!25,
                        minimum size=11mm, align=center, font=\scriptsize,
                        inner sep=0pt},
  op/buffer/.style   = {cylinder, shape border rotate=90, aspect=0.3,
                        draw=arrusBuffer!80!black, thick, fill=arrusBuffer!15,
                        minimum width=18mm, minimum height=10mm, align=center,
                        font=\scriptsize},
  edge/.style        = {-{Latex[length=2mm]}, thick, draw=black!70},
  % Intra-region data flow: shared buffer + CUDA-stream / kernel-order
  % ordering. No runtime queue.
  edge/intra/.style  = {-{Latex[length=2mm]}, thick, draw=black!70},
  % Cross-region stream: real queue infrastructure (backpressure/drop policy,
  % scheduler wake-up). Used for cross-device transfers and rate mismatches.
  edge/stream/.style = {-{Latex[length=2mm]}, very thick, draw=black!70,
                        double distance=1pt},
  edge/trigger/.style= {-{Latex[length=2mm]}, thick, draw=arrusTrigger!80!black,
                        dashed},
  edge/variable/.style = {-{Latex[length=2mm]}, thick,
                          draw=arrusVariable!80!black, densely dotted},
  region/.style      = {rounded corners, inner sep=4mm, dashed,
                        line width=0.8pt},
  region/us4r/.style = {region, draw=arrusUs4R!70!black, fill=arrusUs4R!5},
  region/gpu/.style  = {region, draw=arrusGPU!70!black,  fill=arrusGPU!5},
  legendbox/.style   = {rectangle, draw=black!30, rounded corners=2pt,
                        inner sep=4pt, fill=white},
}

%
% Provides a consistent color palette and node styles for the three figures
% in docs/td/ndarray_design.md.

\usepackage{tikz}
\usetikzlibrary{positioning, arrows.meta, fit, calc, backgrounds, shapes.geometric}

% --- Device / role color palette ------------------------------------------------
\definecolor{arrusUs4R}   {HTML}{2CA9BC}   % ultrasound device (teal)
\definecolor{arrusGPU}    {HTML}{F28C28}   % CUDA GPU (orange)
\definecolor{arrusCPU}    {HTML}{4CAF50}   % CPU (green)
\definecolor{arrusHost}   {HTML}{9E9E9E}   % host / control (gray)
\definecolor{arrusBuffer} {HTML}{616161}   % stream buffer (dark gray)
\definecolor{arrusVariable}{HTML}{F4C542}  % Variable (gold)
\definecolor{arrusTrigger}{HTML}{D64545}   % Trigger source (red)
\definecolor{arrusOutput} {HTML}{5C6BC0}   % Graph output / sink (indigo)

% --- Node styles ---------------------------------------------------------------
\tikzset{
  op/.style       = {rectangle, rounded corners=2pt, draw=black!60, thick,
                     minimum width=26mm, minimum height=8mm, align=center,
                     font=\small},
  op/us4r/.style  = {op, fill=arrusUs4R!25,   draw=arrusUs4R!80!black},
  op/gpu/.style   = {op, fill=arrusGPU!25,    draw=arrusGPU!80!black},
  op/cpu/.style   = {op, fill=arrusCPU!25,    draw=arrusCPU!70!black},
  op/host/.style  = {op, fill=arrusHost!20,   draw=arrusHost!70!black},
  op/output/.style= {op, fill=arrusOutput!25, draw=arrusOutput!80!black},
  op/variable/.style = {ellipse, draw=arrusVariable!70!black, thick,
                        fill=arrusVariable!30, minimum width=22mm,
                        minimum height=7mm, align=center, font=\small\itshape},
  op/trigger/.style  = {regular polygon, regular polygon sides=3,
                        draw=arrusTrigger!80!black, thick, fill=arrusTrigger!25,
                        minimum size=11mm, align=center, font=\scriptsize,
                        inner sep=0pt},
  op/buffer/.style   = {cylinder, shape border rotate=90, aspect=0.3,
                        draw=arrusBuffer!80!black, thick, fill=arrusBuffer!15,
                        minimum width=18mm, minimum height=10mm, align=center,
                        font=\scriptsize},
  edge/.style        = {-{Latex[length=2mm]}, thick, draw=black!70},
  % Intra-region data flow: shared buffer + CUDA-stream / kernel-order
  % ordering. No runtime queue.
  edge/intra/.style  = {-{Latex[length=2mm]}, thick, draw=black!70},
  % Cross-region stream: real queue infrastructure (backpressure/drop policy,
  % scheduler wake-up). Used for cross-device transfers and rate mismatches.
  edge/stream/.style = {-{Latex[length=2mm]}, very thick, draw=black!70,
                        double distance=1pt},
  edge/trigger/.style= {-{Latex[length=2mm]}, thick, draw=arrusTrigger!80!black,
                        dashed},
  edge/variable/.style = {-{Latex[length=2mm]}, thick,
                          draw=arrusVariable!80!black, densely dotted},
  region/.style      = {rounded corners, inner sep=4mm, dashed,
                        line width=0.8pt},
  region/us4r/.style = {region, draw=arrusUs4R!70!black, fill=arrusUs4R!5},
  region/gpu/.style  = {region, draw=arrusGPU!70!black,  fill=arrusGPU!5},
  legendbox/.style   = {rectangle, draw=black!30, rounded corners=2pt,
                        inner sep=4pt, fill=white},
}

%
% Provides a consistent color palette and node styles for the three figures
% in docs/td/ndarray_design.md.

\usepackage{tikz}
\usetikzlibrary{positioning, arrows.meta, fit, calc, backgrounds, shapes.geometric}

% --- Device / role color palette ------------------------------------------------
\definecolor{arrusUs4R}   {HTML}{2CA9BC}   % ultrasound device (teal)
\definecolor{arrusGPU}    {HTML}{F28C28}   % CUDA GPU (orange)
\definecolor{arrusCPU}    {HTML}{4CAF50}   % CPU (green)
\definecolor{arrusHost}   {HTML}{9E9E9E}   % host / control (gray)
\definecolor{arrusBuffer} {HTML}{616161}   % stream buffer (dark gray)
\definecolor{arrusVariable}{HTML}{F4C542}  % Variable (gold)
\definecolor{arrusTrigger}{HTML}{D64545}   % Trigger source (red)
\definecolor{arrusOutput} {HTML}{5C6BC0}   % Graph output / sink (indigo)

% --- Node styles ---------------------------------------------------------------
\tikzset{
  op/.style       = {rectangle, rounded corners=2pt, draw=black!60, thick,
                     minimum width=26mm, minimum height=8mm, align=center,
                     font=\small},
  op/us4r/.style  = {op, fill=arrusUs4R!25,   draw=arrusUs4R!80!black},
  op/gpu/.style   = {op, fill=arrusGPU!25,    draw=arrusGPU!80!black},
  op/cpu/.style   = {op, fill=arrusCPU!25,    draw=arrusCPU!70!black},
  op/host/.style  = {op, fill=arrusHost!20,   draw=arrusHost!70!black},
  op/output/.style= {op, fill=arrusOutput!25, draw=arrusOutput!80!black},
  op/variable/.style = {ellipse, draw=arrusVariable!70!black, thick,
                        fill=arrusVariable!30, minimum width=22mm,
                        minimum height=7mm, align=center, font=\small\itshape},
  op/trigger/.style  = {regular polygon, regular polygon sides=3,
                        draw=arrusTrigger!80!black, thick, fill=arrusTrigger!25,
                        minimum size=11mm, align=center, font=\scriptsize,
                        inner sep=0pt},
  op/buffer/.style   = {cylinder, shape border rotate=90, aspect=0.3,
                        draw=arrusBuffer!80!black, thick, fill=arrusBuffer!15,
                        minimum width=18mm, minimum height=10mm, align=center,
                        font=\scriptsize},
  edge/.style        = {-{Latex[length=2mm]}, thick, draw=black!70},
  % Intra-region data flow: shared buffer + CUDA-stream / kernel-order
  % ordering. No runtime queue.
  edge/intra/.style  = {-{Latex[length=2mm]}, thick, draw=black!70},
  % Cross-region stream: real queue infrastructure (backpressure/drop policy,
  % scheduler wake-up). Used for cross-device transfers and rate mismatches.
  edge/stream/.style = {-{Latex[length=2mm]}, very thick, draw=black!70,
                        double distance=1pt},
  edge/trigger/.style= {-{Latex[length=2mm]}, thick, draw=arrusTrigger!80!black,
                        dashed},
  edge/variable/.style = {-{Latex[length=2mm]}, thick,
                          draw=arrusVariable!80!black, densely dotted},
  region/.style      = {rounded corners, inner sep=4mm, dashed,
                        line width=0.8pt},
  region/us4r/.style = {region, draw=arrusUs4R!70!black, fill=arrusUs4R!5},
  region/gpu/.style  = {region, draw=arrusGPU!70!black,  fill=arrusGPU!5},
  legendbox/.style   = {rectangle, draw=black!30, rounded corners=2pt,
                        inner sep=4pt, fill=white},
}


\begin{document}
\begin{tikzpicture}[
    node distance=5mm and 10mm,
    every node/.append style={font=\small},
    op/arrus/.style = {op, fill=arrusGPU!20, draw=arrusGPU!80!black},
    op/us4rlow/.style = {op, fill=arrusUs4R!25, draw=arrusUs4R!80!black},
]

  % --- Variable input: sos --------------------------------------------------
  \node[op/variable, minimum width=44mm] (sos)
      {Variable(``sos'')\\\scriptsize\itshape NdArrayType((), f32, @ANY)};

  % --- Level 1 sub-ops (arrus dialect, inside arrus.tx_rx) ------------------
  \node[op/arrus, below=18mm of sos, minimum width=88mm] (cd) {%
      \textbf{arrus.compute\_delays}\\%
      \scriptsize\itshape (sos, focus, angle, aperture, geometry)
      $\to$ delays: NdArrayType((n\_el,), f32, @host)};

  \node[op/arrus, below=of cd, minimum width=88mm] (ct) {%
      \textbf{arrus.compose\_tx}\\%
      \scriptsize\itshape (aperture, pulse, delays)
      $\to$ tx\_desc};

  % --- Nested Level 2 body of arrus.acquire (us4r dialect) ------------------
  % These live INSIDE the arrus.acquire region, drawn a bit indented.
  \node[op/us4rlow, below=22mm of ct, minimum width=82mm] (d2c) {%
      \textbf{us4r.delays\_to\_cycles}\\%
      \scriptsize\itshape (delays: f32, clock\_freq)
      $\to$ cycles: NdArrayType((n\_el,), u32, @Us4R:0)};

  \node[op/us4rlow, below=of d2c, minimum width=82mm] (sdr) {%
      \textbf{us4r.set\_delays\_register}\\%
      \scriptsize\itshape (cycles, tx\_id)
      \; [side-effect: writes device DDR shadow buffer]};

  \node[op/us4rlow, below=of sdr, minimum width=82mm] (sam) {%
      \textbf{us4r.set\_aperture\_mask}\\%
      \scriptsize\itshape (aperture\_bits)
      \; [side-effect]};

  \node[op/us4rlow, below=of sam, minimum width=82mm] (at) {%
      \textbf{us4r.arm\_trigger}\\%
      \scriptsize\itshape (trigger\_source, prf)
      \; [side-effect]};

  \node[op/us4rlow, below=of at, minimum width=82mm] (fire) {%
      \textbf{us4r.fire}\\%
      \scriptsize\itshape ()
      $\to$ raw\_frame: NdArrayType((n\_ch, n\_samples), i16, @Us4R:0)};

  % --- Output ---------------------------------------------------------------
  \node[op/output, below=8mm of fire, minimum width=48mm] (raw)
      {raw\_frame};

  % --- Regions (drawn before edges so they sit on background layer) ---------
  \begin{scope}[on background layer]
    % Inner region: arrus.acquire's body (Level-2 us4r dialect).
    % Label placed inside the region at top-left so it doesn't clash with
    % the outer region's label and doesn't run off the standalone canvas.
    \node[region/us4r, fit=(d2c) (sdr) (sam) (at) (fire),
          inner sep=5mm, line width=1pt,
          label={[font=\scriptsize\itshape, text=arrusUs4R!80!black,
                  anchor=south east, xshift=-2mm, yshift=-1mm]
                  north east:\textbf{arrus.acquire} body \textit{(lowered to us4r dialect)}}]
        (acquire_region) {};

    % Outer region: arrus.tx_rx (composite Level-1 op containing three sub-ops).
    % Label placed centered above the region.
    \node[region/gpu, fit=(cd) (ct) (acquire_region),
          inner sep=6mm, line width=1pt,
          label={[font=\small\bfseries, text=arrusGPU!80!black]above:%
                  arrus.tx\_rx \; \textit{(Level 1 --- arrus dialect, composite op)}}]
        (txrx_region) {};
  \end{scope}

  % --- Data-flow edges ------------------------------------------------------
  \draw[edge/variable] (sos.south) -- (cd.north)
      node[midway, right, font=\scriptsize\itshape, text=arrusVariable!70!black]
          {reads current sos};

  \draw[edge/intra] (cd.south) -- (ct.north)
      node[midway, right, font=\scriptsize\itshape, text=black!55]
          {delays};

  % Lowering barrier between ct and d2c.
  \draw[very thick, -{Latex[length=3mm]}, draw=arrusUs4R!80!black]
      (ct.south) -- ++(0,-6mm)
      node[midway, right=1mm, font=\scriptsize\bfseries,
           text=arrusUs4R!80!black]
          {compiled by Us4RBackend.compile}
      -- (d2c.north);

  \draw[edge/intra] (d2c.south) -- (sdr.north)
      node[midway, right, font=\scriptsize\itshape, text=black!55]
          {cycles};

  % These three ops are all side-effect sinks writing to device state;
  % dashed thin arrows show the sequence, not a data flow.
  \draw[-{Latex[length=2mm]}, draw=arrusUs4R!60, dashed, thick]
      (sdr.south) -- (sam.north);
  \draw[-{Latex[length=2mm]}, draw=arrusUs4R!60, dashed, thick]
      (sam.south) -- (at.north);

  % After arm, fire produces the raw frame.
  \draw[edge/intra] (at.south) -- (fire.north);

  \draw[edge/stream] (fire.south) -- (raw.north)
      node[midway, right, font=\scriptsize\itshape, text=black!55]
          {raw echoes};

  % --- Legend ---------------------------------------------------------------
  \node[legendbox, below=14mm of raw, anchor=north] {%
    \begin{tikzpicture}[node distance=3mm]
      \node[op/arrus,    minimum width=26mm, minimum height=6mm] (l1) {arrus.* (Level 1)};
      \node[op/us4rlow,  minimum width=26mm, minimum height=6mm, below=of l1] (l2) {us4r.* (Level 2)};
      \node[op/variable, minimum width=26mm, minimum height=6mm, below=of l2] (l3) {Variable};
      \node[op/output,   minimum width=26mm, minimum height=6mm, below=of l3] (l4) {output};
      \node[below=3mm of l4.south, anchor=north, font=\scriptsize] (e1)
        {\tikz\draw[edge/intra]  (0,0) -- (8mm,0);\ SSA value};
      \node[below=1mm of e1.south, anchor=north, font=\scriptsize] (e2)
        {\tikz\draw[edge/variable] (0,0) -- (8mm,0);\ variable read};
      \node[below=1mm of e2.south, anchor=north, font=\scriptsize] (e3)
        {\tikz\draw[very thick, -{Latex[length=2mm]}, draw=arrusUs4R!80!black]
              (0,0) -- (8mm,0);\ lowering};
      \node[below=1mm of e3.south, anchor=north, font=\scriptsize] (e4)
        {\tikz\draw[-{Latex[length=2mm]}, draw=arrusUs4R!60, dashed, thick]
              (0,0) -- (8mm,0);\ sequenced side-effects};
    \end{tikzpicture}%
  };

\end{tikzpicture}
\end{document}
